\section{Load sequence files through a reusable parser}
\label{sec:arch-ingestion}
\label{sec:io-loading-contract}

Public sequence archives distribute FASTA or FASTQ text, often compressed with
gzip. BioSeqDB's input layer turns that text into compact sequence chunks and
raw headers. A PostgreSQL sink then turns the headers into typed columns and
stores the chunks. Separating these jobs lets the same parser serve database
loading and ordinary C++ consumers without importing PostgreSQL into the
file-reading code.

\subsection{Separate byte sources from record framing}
\label{sec:arch-io-sources}
\label{sec:io-reader-architecture}
\label{sec:io-parsing}

A byte source fills a caller-owned buffer and reports bytes read, end of input,
or an error. The parser is specialized for the source type, so refill calls
need no virtual dispatch. Plain files can use a memory mapping or sequential
reads. The SQL loading path uses zlib for gzip input. Biological format and
compression are selected independently: the filename suffix chooses FASTA or
FASTQ, while gzip magic bytes identify compression. The loader does not infer
the biological grammar from an arbitrary first record.

The parser design was inspired by Helicase's vectorized parsing and compact
DNA output, and by StringZilla's vectorized byte-search
operations~\autocite{helicase,stringzilla}. BioSeqDB implements its own kernels
through Highway. These libraries are not runtime dependencies of the input
path. In FASTA, the header begins with \texttt{>}, followed by a sequence body.
For DNA4 packed into 64-bit units, the optimized body loop searches for the
next \texttt{>} while encoding symbols. Boundary detection and packing thus
share an input pass. Other alphabet and unit combinations use the same record
states with their corresponding encoder and line scanner.

FASTQ adds a \texttt{+} separator and quality bytes. Quality lines can contain
\texttt{@}, \texttt{+}, or \texttt{>}, so those characters cannot independently
identify the next record. The parser instead counts quality bytes against the
raw sequence length after removing line endings. This count is distinct from
the number of symbols retained by the compact encoder.

\subsection{Deliver chunks while a record remains open}
\label{sec:arch-io-chunks}
\label{sec:io-bounds}

The parser invokes sink functions, called callbacks, when a record begins,
when packed chunks are available, and when the record ends. It keeps its
current state and an unfinished
storage word across refills. Complete words can be delivered before the record
ends. The final chunk supplies the remaining word and its valid symbol count.
Figure~\ref{fig:arch-io-state} shows why an input-buffer boundary need not be a
record boundary.

\begin{figure}[htbp]
\centering
\begin{tikzpicture}[x=1cm,y=1cm,font=\footnotesize,
  source/.style={draw=black!40,fill=black!2},
  record a/.style={draw=blue!65!black,fill=blue!9},
  record b/.style={draw=green!45!black,fill=green!10},
  flow/.style={-{Latex[length=2mm]},thick,draw=black!65}]
  \node[anchor=west,font=\footnotesize\bfseries] at (.25,6.4)
    {Source bytes: record A has 70 bases, record B has 40};
  \foreach \left/\right/\middle/\name in
    {.25/4.45/2.35/1,4.65/8.85/6.75/2,9.05/13.25/11.15/3} {
    \draw[source] (\left,3.95) rectangle (\right,5.85);
    \node at (\middle,6.05) {input buffer \name};
  }
  \draw[record a] (.45,4.7) rectangle (4.25,5.6);
  \node[align=center] at (2.35,5.15)
    {\texttt{>A} + LF\\A: bases 1--38};
  \draw[record a] (4.85,5.0) rectangle (8.65,5.6);
  \node[align=center] at (6.75,5.3) {A: bases 39--70 + LF};
  \draw[record b] (4.85,4.05) rectangle (8.65,4.85);
  \node[align=center] at (6.75,4.45)
    {\texttt{>B} + LF\\B: bases 1--5};
  \draw[record b] (9.25,4.15) rectangle (13.05,5.05);
  \node[align=center] at (11.15,4.6)
    {B: bases 6--40 + LF\\end of file};

  \foreach \middle in {2.35,6.75,11.15}
    \draw[flow] (\middle,3.95) -- (\middle,3.2);
  \node at (2.35,3.0) {A: first chunk};
  \node at (6.75,3.0) {A: final chunk};
  \node at (11.15,3.0) {B: final chunk};

  \draw[record a] (1.65,2.1) rectangle (3.05,2.7);
  \node at (2.35,2.4) {32};
  \draw[record a] (5.35,2.1) rectangle (6.75,2.7);
  \draw[record a,fill=blue!3,dashed] (6.75,2.1) rectangle (8.15,2.7);
  \node at (6.05,2.4) {32};
  \node at (7.45,2.4) {6 + pad};
  \draw[record b] (9.75,2.1) rectangle (11.15,2.7);
  \draw[record b,fill=green!3,dashed] (11.15,2.1) rectangle (12.55,2.7);
  \node at (10.45,2.4) {32};
  \node at (11.85,2.4) {8 + pad};
  \node[align=center] at (2.35,1.55)
    {6 bases remain\\in A's partial word};
  \node[align=center] at (6.75,1.55)
    {A closes at 70 bases\\B retains 5 bases};
  \node[align=center] at (11.15,1.55)
    {B closes at 40 bases\\with 8 final valid bases};
  \node[anchor=west,font=\footnotesize\bfseries] at (.25,.65)
    {Packed output: each box is one 64-bit DNA4 word};
\end{tikzpicture}
\caption{Two single-line FASTA records parsed with illustrative 41-byte input
buffers. LF denotes a line-feed byte. Header and LF bytes occupy input space
but are not sequence symbols.
The first refill ends inside A: 32 bases are delivered, while six remain in
its partial word. The next buffer finishes A and starts B. Source and packed
widths are schematic. Word labels count valid bases, with padding in dashed
words. Chunk views
are borrowed for their callback and copied by the backend before reuse.}
\label{fig:arch-io-state}
\label{fig:io-load-pipeline}
\end{figure}

The SQL sink runs on the PostgreSQL backend thread. Its default path consumes
callbacks directly. Optional overlap moves owned events through a FIFO queue
of fixed capacity: a worker parses, and the backend performs database writes.
A full queue blocks the producer. Closing it wakes blocked participants, and
the worker is joined before a consumer error is propagated.

The reusable batch API has a different contract. It owns complete batches and
can use direct mapped views, parallel parser shards, or a parallel gzip
source. Those options do not describe SQL streaming, which refills a slab and
uses the single-threaded zlib source for gzip. In particular, a batch parser's
throughput is not a measurement of database loading.

\subsection{Route header fields into SQL columns}
\label{sec:arch-io-headers}
\label{sec:io-header-mapping}
\label{sec:io-sql}

Headers can be preserved verbatim or parsed using a named convention, such as
GenBank, RefSeq, or UniProt. Their conventions distinguish identifiers from
annotations carried in the same definition
line~\autocite{ncbiFastaFormat,uniprotFastaHeaders}. Automatic selection samples
a bounded decompressed prefix once per file. The loader also resolves field-to-column routes once,
using destination names and any explicit mappings. Per-record parsing then
passes field slices to the PostgreSQL sink, which applies the destination
types. A header dialect describes metadata structure. It is not a strict
validator for the sequence alphabet.

For example, a gzipped FASTA file may contain the header
\texttt{>dna\_new [organism=Homo sapiens] [chromosome=17] [moltype=mRNA]}.
The following table gives those fields their relational types:

\noindent\begin{minipage}{\linewidth}
\begin{lstlisting}[style=bioseq-arch-sql,basicstyle=\ttfamily\footnotesize,
  caption={Load a FASTA file with GenBank-style header fields.},
  label={lst:arch-load-genbank}]
CREATE TABLE sequences (
  id text,
  organism text,
  chromosome integer,
  moltype text,
  seq bioseq.seq('dna4', 64)
) USING bioseq;

SELECT bioseq.load('sequences'::regclass,
                   '/data/example.fasta.gz');
\end{lstlisting}
\end{minipage}

Here automatic routing produces identifier \texttt{dna\_new}, organism
\texttt{Homo sapiens}, integer chromosome \texttt{17}, and molecule type
\texttt{mRNA}. The return value is the number of loaded records. Compression
changes the byte source, while the SQL interface remains the same.

The sequence column selects the compact alphabet and storage-unit width
before parsing starts. Once a record's fragments are ready, the sink writes
the parent row and publishes its sequence locator. PostgreSQL controls
transaction visibility and rollback. The sink evaluates defaults and
\texttt{NOT NULL}/\texttt{CHECK} constraints. It does not run ordinary
\texttt{INSERT} triggers or maintain ordinary secondary indexes. Normal
trigger-based foreign-key checks and index-based unique or exclusion checks
therefore do not run.

\label{sec:arch-io-limits}
\label{sec:io-integrity-limits}
The memory guarantee concerns streamed packed bodies. Headers and FASTQ
quality storage can grow with record length, and paths that retain unfinished
physical lines remain subject to the input-buffer limit. Alphabet handling
is permissive, and the declared strict-alphabet option is not enforced by the
parser. FASTQ length checks also have an exception for excess quality bytes
on an unterminated final line. These limits preclude treating successful
loading as complete input validation.
